% ==============================================================================
%  1.3  Tipos de datos según su procedencia o presentación
% ==============================================================================
\section{Tipos de datos según su procedencia}
\label{sec:t1-datos}

\subsection{Datos observacionales y experimentales}

Según procedan de estudios no experimentales o experimentales, los datos se clasifican en
\emph{observacionales} y \emph{experimentales}.

\begin{definicion}[Datos observacionales y experimentales]
  \begin{itemize}
    \item \textbf{Datos observacionales}: se obtienen de estudios no experimentales, en los
      que \emph{no se controlan} las variables explicativas de interés. Esto limita las
      conclusiones, ya que es muy difícil establecer relaciones causa--efecto.
    \item \textbf{Datos experimentales}: se obtienen de experimentos controlados y permiten
      extraer conclusiones más sólidas sobre posibles relaciones causa--efecto. La
      \emph{aleatorización} es un mecanismo importante para controlar las variables
      explicativas.
  \end{itemize}
\end{definicion}

Un experimento controlado tiene como objetivo comparar el efecto sobre la respuesta de
distintos \textbf{tratamientos} (combinaciones de niveles de las variables explicativas),
aplicándolos a diferentes \textbf{unidades experimentales} (individuos u objetos sobre los que
se realiza la investigación).

\begin{definicion}[Aleatorización]
  La \textbf{aleatorización} consiste en asignar de forma aleatoria cada unidad
  experimental a un tratamiento. Es útil porque tiende a equilibrar los posibles efectos de
  otras variables (no controladas) que podrían afectar a la respuesta.
\end{definicion}

\begin{ejemplo}[Efecto de fumar en la capacidad pulmonar]
  \begin{center}
    \small
    \begin{tabular}{p{0.45\linewidth}p{0.45\linewidth}}
      \toprule
      \textbf{Diseño experimental} & \textbf{Diseño observacional} \\
      \midrule
      Elegir 100 personas que actualmente no fumen. & --- \\[2pt]
      Asignar \textbf{aleatoriamente} 50 de ellas al grupo ``fumadores'' y las otras 50 al
      grupo ``no fumadores''. & --- \\[2pt]
      Cada persona del grupo ``fumadores'' fumará un paquete al día durante 10 años. &
      Encontrar a 50 personas fumadoras de un paquete al día durante 10 años. \\[2pt]
      Cada persona del grupo ``no fumadores'' permanecerá sin fumar durante 10 años. &
      Encontrar a 50 personas que no hayan fumado durante 10 años. \\[2pt]
      Transcurridos los 10 años, medir la capacidad pulmonar de las 100 personas. &
      Medir la capacidad pulmonar de las 100 personas. \\
      \bottomrule
    \end{tabular}
  \end{center}
  En el diseño observacional, los fumadores y los no fumadores pueden diferir en otras
  variables (edad, trabajo, deporte\ldots) que afectan a la capacidad pulmonar; en el
  experimental, la aleatorización equilibra esos efectos entre los dos grupos.
\end{ejemplo}

\paragraph{Estudios muestrales.} Entre ambos tipos de estudio se sitúan los \textbf{estudios
muestrales}: las inferencias se realizan sobre la población adecuada, pero la presencia de
variables influyentes no controladas o desconocidas puede dificultar la interpretación de las
relaciones entre las variables observadas.

\paragraph{Precauciones en el análisis.} Al analizar los datos hay que tener en cuenta:
\begin{itemize}
  \item posibles \textbf{cambios en las condiciones de recogida} de los datos (de personal, de
    máquinas, envejecimiento\ldots) que puedan influir en los resultados;
  \item la presencia de variables explicativas \textbf{fuertemente correlacionadas} entre sí,
    cuyos efectos se confunden;
  \item un \textbf{rango pequeño} de alguna variable importante, que impide establecer su
    influencia fuera de ese rango.
\end{itemize}

\subsection{Efectos fijos, aleatorios y mixtos}

Según cómo se recojan las variables explicativas se distingue entre:

\begin{definicion}[Modelos de efectos fijos, aleatorios y mixtos]
  \begin{itemize}
    \item \textbf{Modelo de efectos fijos.} Una \emph{variable fija} es aquella que se supone
      medida sin error, de modo que los valores observados en un estudio serían los mismos en
      otro. Las variables explicativas toman valores predeterminados, que se controlan en el
      diseño y la realización del experimento.
    \item \textbf{Modelo de efectos aleatorios.} Las variables explicativas son \emph{variables
      aleatorias}, cuyos valores son una muestra aleatoria de todos los valores posibles, y se
      espera poder generalizar los resultados al resto de valores. Esta generalización debe
      hacerse con cautela, ya que los resultados solo son válidos en el rango de variación
      conjunta de las variables implicadas.
    \item \textbf{Modelo mixto}: combina efectos fijos y aleatorios. En muchas aplicaciones los
      modelos son mixtos.
  \end{itemize}
\end{definicion}

\begin{ejemplo}[Dosis de un fármaco]
  Se mide la efectividad de un fármaco con dosis de 5, 10 y 15\,mg.
  \begin{itemize}
    \item Si esas tres dosis se han elegido intencionadamente y los resultados solo se van a
      aplicar a ellas, ``dosis'' es un efecto \textbf{fijo}.
    \item Si se consideran una muestra aleatoria de todas las dosis posibles entre 5 y 15\,mg
      y los resultados se quieren aplicar a cualquier dosis de ese rango, ``dosis'' es un
      efecto \textbf{aleatorio}.
  \end{itemize}
\end{ejemplo}
